\section{Протокол}
\label{sec:architecture}

\textbf{Память модели.} У Qwen3.8-27B два вида слоёв. Шестнадцать слоёв полного внимания хранят ключи и значения для
каждого прочитанного токена; это KV-кэш, он растёт с контекстом. Остальные слои используют линейное внимание и
хранят одну сводку всей истории фиксированного размера, которая обновляется с каждым токеном; это рекуррентное
состояние. Мы измерили обе части на сохранённом слоте llama.cpp: KV-кэш в 8 битах занимает 34\,844 байта на токен,
рекуррентное состояние 156.9~МБ при любой длине контекста, потому что движок хранит его в 32-битных числах. Веса в 4
битах занимают около 19~ГБ, и на машине с 64~ГБ рядом с ними помещаются три сессии по 256 тысяч токенов. Дельтой мы
называем то, что сервер отправляет клиенту за один вызов: KV-кэш новых позиций и рекуррентное состояние целиком. Это
и есть массив чисел, который копирует схема.

\textbf{Ход с оффлоадом.} Рисунок~\ref{fig:sequence} показывает два вида хода. Клиент отправляет серверу
идентификаторы токенов, которые появились после прошлой синхронизации: вывод инструмента и собственный ответ модели
на прошлом ходе. Сервер дописывает их к своей копии сессии, считает prefill для всех новых позиций, кроме последней, и
возвращает дельту. Клиент вставляет дельту в свой кэш и сам считает последний токен порции, потому что вставленная
память не несёт распределения следующего токена. Затем клиент генерирует ответ. Вверх идут килобайты, вниз мегабайты:
отправлять кэш вверх слишком дорого, поэтому копия сессии у сервера пополняется только идентификаторами, а
недостающие позиции сервер досчитывает сам.

\begin{figure}[t]
\centering
\begin{tikzpicture}[x=1mm,y=1mm]
\node[font=\footnotesize\bfseries] at (0,0) {Клиент};
\node[font=\footnotesize\bfseries] at (60,0) {Сервер};
\draw[thick] (0,-3) -- (0,-62);
\draw[thick] (60,-3) -- (60,-62);
\node[lbl,anchor=west] at (-2,-8) {\textbf{Ход локально:} малая порция};
\draw[thick] (-1.5,-12) rectangle (1.5,-20);
\node[lbl,anchor=west] at (3,-16) {prefill порции и decode на клиенте};
\draw[arr,dashed] (1.5,-23) -- node[lbl,above]{id токенов (фоном)} (58.5,-23);
\node[lbl,anchor=west] at (-2,-30) {\textbf{Ход с оффлоадом:} крупная порция};
\draw[arr] (1.5,-35) -- node[lbl,above]{id новых токенов} (58.5,-35);
\draw[thick] (58.5,-37) rectangle (61.5,-45);
\node[lbl,anchor=west] at (63,-41) {prefill новых позиций};
\draw[arr] (58.5,-48) -- node[lbl,above]{дельта: KV и состояние} (1.5,-48);
\draw[thick] (-1.5,-50) rectangle (1.5,-60);
\node[lbl,anchor=west] at (3,-55) {вставка дельты, последний токен, decode};
\end{tikzpicture}
\caption{Два вида хода. Малую порцию клиент читает сам и отправляет серверу только идентификаторы, чтобы копия сессии
у сервера оставалась полной. Крупную порцию читает сервер и возвращает дельту; клиент считает только последний токен
порции и генерирует ответ.}
\label{fig:sequence}
\end{figure}

\textbf{Сервер хранит сессию.} Сервер держит у себя KV-кэш и рекуррентное состояние каждой активной сессии и считает
только приращение. Если бы сервер не помнил сессию, каждый вызов требовал бы чтения всего контекста заново: на средней
глубине профиля, 164 тысячи токенов, это в 100--160 раз больше работы, чем чтение тысячи новых токенов.

\textbf{Сессии только растут.} Клиент гибрида работает как локальный прокси между агентной средой и движком
генерации и отвечает среде по тому же API, что и обычный сервер модели, поэтому схема работает с любой агентной
средой, у которой история разговора только накапливается. Клиент ведёт журнал сессии: точные идентификаторы токенов,
которые лежат в его кэше, и границы сообщений. Новый запрос агента клиент сверяет с журналом и дописывает только новые
сообщения. Если среда сжимает старый контекст, кэш сбрасывается в любой системе, облачной и локальной: клиент
откатывается по журналу к последней контрольной точке до расхождения, и всё, что после неё, сервер читает как крупную
порцию. Если сервер недоступен, клиент читает порции сам и только ждёт дольше.

\textbf{Правило оффлоада.} Обозначим через $d$ число новых токенов в вызове, через $L$ глубину контекста, через $B$
скорость канала. Для $\ell$ слоёв полного внимания с $h$ головами ключей и значений размерности $e$ и $b$ байтами на
число кэш на токен равен $K = 2 \ell h e b$, и дельта весит $D(d) = d K + S$, где $S$ размер рекуррентного состояния.
Вычисления на токен prefill $F(L) \approx 2N + 4 \ell L q e$, где $N$ число активных параметров, $q$ число голов
запросов; второе слагаемое растёт с глубиной. Локальное чтение занимает $t_{\text{loc}} = d / r(L)$, где $r(L)$
замеренная скорость prefill клиента на глубине $L$. Ход с оффлоадом занимает $t_{\text{off}} = T_0 + t_{\text{up}} +
d F(L) / (P u) + D(d) / B + t_{\text{ins}} + t_{\text{last}}$, где $T_0$ постоянная задержка вызова, $t_{\text{up}}$
подъём копии сессии в память GPU, $P$ пиковая производительность GPU сервера, $u$ доля пика, которую использует
движок, $t_{\text{ins}}$ вставка дельты и $t_{\text{last}}$ чтение последнего токена на клиенте. Клиент отдаёт порцию
серверу, если $t_{\text{off}} < t_{\text{loc}}$; очередей движок расчётов не моделирует (раздел~\ref{sec:limits}).
